\documentclass[12pt]{article}

% Idioma y codificación
\usepackage[spanish, es-tabla, es-notilde]{babel}       %es-tabla para que se titule "Tabla"
\usepackage[utf8]{inputenc}

% Márgenes
\usepackage[a4paper,top=3cm,bottom=2.5cm,left=3cm,right=3cm]{geometry}

% Comentarios de bloque
\usepackage{verbatim}

% Paquetes de links
\usepackage[hidelinks]{hyperref}    % Permite enlaces
\usepackage{url}                    % redirecciona a la web

% Más opciones para enumeraciones
\usepackage{enumitem}

% Personalizar la portada
\usepackage{titling}

% Paquetes de tablas
\usepackage{multirow}

% Para añadir el símbolo de euro
\usepackage{eurosym}

% Paquete para subrayar con diferentes estilos
\usepackage[normalem]{ulem}


%------------------------------------------------------------------------

%Paquetes de figuras
\usepackage{caption}
\usepackage{subcaption} % Figuras al lado de otras
\usepackage{float}      % Poner figuras en el sitio indicado H.


% Paquetes de imágenes
\usepackage{graphicx}       % Paquete para añadir imágenes
\usepackage{transparent}    % Para manejar la opacidad de las figuras

% Paquete para usar colores
\usepackage[dvipsnames, table, xcdraw]{xcolor}
\usepackage{pagecolor}      % Para cambiar el color de la página

% Habilita tamaños de fuente mayores
\usepackage{fix-cm}

% Para los gráficos
\usepackage{tikz}
\usepackage{forest}

% Para poder situar los nodos en los grafos
\usetikzlibrary{positioning}

% Para evitar problemas entre babel y tikz
\usetikzlibrary{babel}


%------------------------------------------------------------------------

% Paquetes de matemáticas
\usepackage{mathtools, amsfonts, amssymb, mathrsfs}
\usepackage[makeroom]{cancel}     % Simplificar tachando
\usepackage{polynom}    % Divisiones y Ruffini
\usepackage{units} % Para poner fracciones diagonales con \nicefrac

\usepackage{pgfplots}   %Representar funciones
\pgfplotsset{compat=1.18}  % Versión 1.18

\usepackage{tikz-cd}    % Para usar diagramas de composiciones
\usetikzlibrary{calc}   % Para usar cálculo de coordenadas en tikz

%Definición de teoremas, etc.
\usepackage{amsthm}
%\swapnumbers   % Intercambia la posición del texto y de la numeración

\theoremstyle{plain}

\makeatletter
\@ifclassloaded{article}{
  % --- Configuración para ARTICLE ---
  \newtheorem{teo}{Teorema}[section]
  \theoremstyle{definition}
  \newtheorem{definicion}{Definición}[section]
  \newtheorem{ejercicio}{Ejercicio}[section]

  \renewcommand{\theteo}{%
    \ifnum\value{section}=0 \arabic{teo}\else \thesection.\arabic{teo}\fi
  }
  \renewcommand{\thedefinicion}{%
    \ifnum\value{section}=0 \arabic{definicion}\else \thesection.\arabic{definicion}\fi
  }
  \renewcommand{\theejercicio}{%
    \ifnum\value{section}=0 \arabic{ejercicio}\else \thesection.\arabic{ejercicio}\fi
  }
}{
  % --- Configuración para BOOK / REPORT ---
  \newtheorem{teo}{Teorema}[chapter]
  \theoremstyle{definition}
  \newtheorem{definicion}{Definición}[chapter]
  \newtheorem{ejercicio}{Ejercicio}[section]

  % Si no hay sección: Capítulo.Teorema (ej. 2.1)
  % Si hay sección:    Capítulo.Sección.Teorema (ej. 2.3.1)
  \renewcommand{\theteo}{%
    \ifnum\value{section}=0
      \thechapter.\arabic{teo}%
    \else
      \thesection.\arabic{teo}% (\thesection ya contiene el número de capítulo)
    \fi
  }
  \renewcommand{\thedefinicion}{%
    \ifnum\value{section}=0
      \thechapter.\arabic{definicion}%
    \else
      \thesection.\arabic{definicion}%
    \fi
  }
  \renewcommand{\theejercicio}{%
    \ifnum\value{section}=0
      \arabic{ejercicio}%
    \else
      \thesection.\arabic{ejercicio}% Formato por defecto: ej. 2.2.1
    \fi
  }
}
\makeatother

% Teoremas dependientes o con contador compartido
\newtheorem{coro}{Corolario}[teo]
\newtheorem{prop}[teo]{Proposición}
\newtheorem{lema}[teo]{Lema}
\newtheorem*{lema*}{Lema}

\theoremstyle{remark}
\newtheorem*{observacion}{Observación}

\theoremstyle{definition}
\newtheorem*{notacion}{Notación}
\newtheorem*{ejemplo}{Ejemplo}
\newtheorem*{ejercicio*}{Ejercicio}


% \renewcommand\qedsymbol{$\blacksquare$}         % Cambiar símbolo QED
%------------------------------------------------------------------------

% Paquetes para encabezados
\usepackage{fancyhdr}
\pagestyle{fancy}
\fancyhf{}

\newcommand{\helv}{ % Modificación tamaño de letra
\fontfamily{}\fontsize{12}{12}\selectfont}
\setlength{\headheight}{15pt} % Amplía el tamaño del índice


%\usepackage{lastpage}   % Referenciar última pag   \pageref{LastPage}
%\fancyfoot[C]{%
%  \begin{minipage}{\textwidth}
%    \centering
%    ~\\
%    \thepage\\
%    \href{https://losdeldgiim.github.io/}{\texttt{\footnotesize losdeldgiim.github.io}}
%  \end{minipage}
%}
\fancyfoot[C]{\thepage}
\fancyfoot[R]{\href{https://losdeldgiim.github.io/}{\texttt{\footnotesize losdeldgiim.github.io}}}

%------------------------------------------------------------------------

% Conseguir que no ponga "Capítulo 1". Sino solo "1."
\makeatletter
\@ifclassloaded{book}{
  \renewcommand{\chaptermark}[1]{\markboth{\thechapter.\ #1}{}} % En el encabezado
    
  \renewcommand{\@makechapterhead}[1]{%
  \vspace*{50\p@}%
  {\parindent \z@ \raggedright \normalfont
    \ifnum \c@secnumdepth >\m@ne
      \huge\bfseries \thechapter.\hspace{1em}\ignorespaces
    \fi
    \interlinepenalty\@M
    \Huge \bfseries #1\par\nobreak
    \vskip 40\p@
  }}
}
\makeatother

%------------------------------------------------------------------------
% Paquetes de cógido
\usepackage{minted}
\renewcommand\listingscaption{Código fuente}

\usepackage{fancyvrb}
% Personaliza el tamaño de los números de línea
\renewcommand{\theFancyVerbLine}{\small\arabic{FancyVerbLine}}

% Estilo para C++
\newminted{cpp}{
    frame=lines,
    framesep=2mm,
    baselinestretch=1.2,
    linenos,
    escapeinside=||
}

% para minted
\definecolor{LightGray}{rgb}{0.95,0.95,0.92}
\setminted{
    linenos=true,
    stepnumber=5,
    numberfirstline=true,
    autogobble,
    breaklines=true,
    breakautoindent=true,
    breaksymbolleft=,
    breaksymbolright=,
    breaksymbolindentleft=0pt,
    breaksymbolindentright=0pt,
    breaksymbolsepleft=0pt,
    breaksymbolsepright=0pt,
    fontsize=\footnotesize,
    bgcolor=LightGray,
    numbersep=10pt
}


\usepackage{listings} % Para incluir código desde un archivo

\renewcommand\lstlistingname{Código Fuente}
\renewcommand\lstlistlistingname{Índice de Códigos Fuente}

% Definir colores
\definecolor{vscodepurple}{rgb}{0.5,0,0.5}
\definecolor{vscodeblue}{rgb}{0,0,0.8}
\definecolor{vscodegreen}{rgb}{0,0.5,0}
\definecolor{vscodegray}{rgb}{0.5,0.5,0.5}
\definecolor{vscodebackground}{rgb}{0.97,0.97,0.97}
\definecolor{vscodelightgray}{rgb}{0.9,0.9,0.9}

% Configuración para el estilo de C similar a VSCode
\lstdefinestyle{vscode_C}{
  backgroundcolor=\color{vscodebackground},
  commentstyle=\color{vscodegreen},
  keywordstyle=\color{vscodeblue},
  numberstyle=\tiny\color{vscodegray},
  stringstyle=\color{vscodepurple},
  basicstyle=\scriptsize\ttfamily,
  breakatwhitespace=false,
  breaklines=true,
  captionpos=b,
  keepspaces=true,
  numbers=left,
  numbersep=5pt,
  showspaces=false,
  showstringspaces=false,
  showtabs=false,
  tabsize=2,
  frame=tb,
  framerule=0pt,
  aboveskip=10pt,
  belowskip=10pt,
  xleftmargin=10pt,
  xrightmargin=10pt,
  framexleftmargin=10pt,
  framexrightmargin=10pt,
  framesep=0pt,
  rulecolor=\color{vscodelightgray},
  backgroundcolor=\color{vscodebackground},
}

%------------------------------------------------------------------------
% La etiqueda IA sirve para indicar a partir de donde le pasamos el preambulo a la IA para el escaneo
% IA

% Comandos definidos
\newcommand{\bb}[1]{\mathbb{#1}}
\newcommand{\cc}[1]{\mathcal{#1}}

% I prefer the slanted \leq
\let\oldleq\leq % save them in case they're every wanted
\let\oldgeq\geq
\renewcommand{\leq}{\leqslant}
\renewcommand{\geq}{\geqslant}

% Si y solo si
\newcommand{\sii}{\iff}

% MCD y MCM
\DeclareMathOperator{\mcd}{mcd}
\DeclareMathOperator{\mcm}{mcm}

% Signo
\DeclareMathOperator{\sgn}{sgn}

% Imagen
\DeclareMathOperator{\Ima}{Im}

% Aplicación identidad
\newcommand{\id}{\mathrm{Id}}

% Letras griegas
\newcommand{\eps}{\epsilon}
\newcommand{\veps}{\varepsilon}
\newcommand{\lm}{\lambda}

\newcommand{\ol}{\overline}
\newcommand{\ul}{\underline}
\newcommand{\wt}{\widetilde}
\newcommand{\wh}{\widehat}

\let\oldvec\vec
\renewcommand{\vec}{\overrightarrow}

% Derivadas parciales
\newcommand{\del}[2]{\frac{\partial #1}{\partial #2}}
\newcommand{\Del}[3]{\frac{\partial^{#1} #2}{\partial #3^{#1}}}
\newcommand{\deld}[2]{\dfrac{\partial #1}{\partial #2}}
\newcommand{\Deld}[3]{\dfrac{\partial^{#1} #2}{\partial #3^{#1}}}


\newcommand{\AstIg}{\stackrel{(\ast)}{=}}
\newcommand{\Hop}{\stackrel{L'H\hat{o}pital}{=}}

\newcommand{\red}[1]{{\color{red}#1}} % Para integrales, destacar los cambios.

% Método de integración
\newcommand{\MetInt}[2]{
    \left[\begin{array}{c}
        #1 \\ #2
    \end{array}\right]
}

% Declarar aplicaciones
% 1. Nombre aplicación
% 2. Dominio
% 3. Codominio
% 4. Variable
% 5. Imagen de la variable
\newcommand{\Func}[5]{
    \begin{equation*}
        \begin{array}{rrll}
            \displaystyle #1:& \displaystyle  #2 & \longrightarrow & \displaystyle  #3\\
               & \displaystyle  #4 & \longmapsto & \displaystyle  #5
        \end{array}
    \end{equation*}
}

%------------------------------------------------------------------------


\newcommand{\myhref}[2]{\href{#1}{\color{blue}\uline{#2}}}

\usetikzlibrary{arrows.meta, decorations.markings} % Cargar las bibliotecas necesarias
% Configuración para las flechas
\tikzset{
    arrow at 1/3/.style={postaction={decorate},
        decoration={markings, mark=at position 0.33 with {\arrow{Stealth}}}},
    arrow at 2/3/.style={postaction={decorate},
        decoration={markings, mark=at position 0.66 with {\arrow{Stealth}}}}
}

\begin{document}

    % 1. Foto de fondo
    % 2. Título
    % 3. Encabezado Izquierdo
    % 4. Color de fondo
    % 5. Coord x del titulo
    % 6. Coord y del titulo
    % 7. Fecha

    
    % 1. Foto de fondo
% 2. Título
% 3. Encabezado Izquierdo
% 4. Color de fondo
% 5. Coord x del titulo
% 6. Coord y del titulo
% 7. Fecha
% 8. Autor

\newcommand{\portada}[8]{

    \portadaBase{#1}{#2}{#3}{#4}{#5}{#6}{#7}{#8}
    \portadaBook{#1}{#2}{#3}{#4}{#5}{#6}{#7}{#8}
}

\newcommand{\portadaExamen}[8]{

    \portadaBase{#1}{#2}{#3}{#4}{#5}{#6}{#7}{#8}
    \portadaArticle{#1}{#2}{#3}{#4}{#5}{#6}{#7}{#8}
}




\newcommand{\portadaBase}[8]{

    % Tiene la portada principal y la licencia Creative Commons
    
    % 1. Foto de fondo
    % 2. Título
    % 3. Encabezado Izquierdo
    % 4. Color de fondo
    % 5. Coord x del titulo
    % 6. Coord y del titulo
    % 7. Fecha
    % 8. Autor
    
    
    \thispagestyle{empty}               % Sin encabezado ni pie de página
    \newgeometry{margin=0cm}        % Márgenes nulos para la primera página
    
    
    % Encabezado
    \fancyhead[L]{\helv #3}
    \fancyhead[R]{\helv \nouppercase{\leftmark}}
    
    
    \pagecolor{#4}        % Color de fondo para la portada
    
    \begin{figure}[p]
        \centering
        \transparent{0.3}           % Opacidad del 30% para la imagen
        
        \includegraphics[width=\paperwidth, keepaspectratio]{../../_assets/#1}
    
        \begin{tikzpicture}[remember picture, overlay]
            \node[anchor=north west, text=white, opacity=1, font=\fontsize{60}{90}\selectfont\bfseries\sffamily, align=left] at (#5, #6) {#2};
            
            \node[anchor=south east, text=white, opacity=1, font=\fontsize{12}{18}\selectfont\sffamily, align=right] at (9.7, 3) {\href{https://losdeldgiim.github.io/}{\textbf{Los Del DGIIM}, \texttt{\footnotesize losdeldgiim.github.io}}};
            
            \node[anchor=south east, text=white, opacity=1, font=\fontsize{12}{15}\selectfont\sffamily, align=right] at (9.7, 1.8) {Doble Grado en Ingeniería Informática y Matemáticas\\Universidad de Granada};
        \end{tikzpicture}
    \end{figure}
    
    
    \restoregeometry        % Restaurar márgenes normales para las páginas subsiguientes
    \pagecolor{white}       % Restaurar el color de página
    
    
    \newpage
    \thispagestyle{empty}               % Sin encabezado ni pie de página
    \begin{tikzpicture}[remember picture, overlay]
        \node[anchor=south west, inner sep=3cm] at (current page.south west) {
            \begin{minipage}{0.5\paperwidth}
                \href{https://creativecommons.org/licenses/by-nc-nd/4.0/}{
                    \includegraphics[height=2cm]{../../_assets/Licencia.png}
                }\vspace{1cm}\\
                Esta obra está bajo una
                \href{https://creativecommons.org/licenses/by-nc-nd/4.0/}{
                    Licencia Creative Commons Atribución-NoComercial-SinDerivadas 4.0 Internacional (CC BY-NC-ND 4.0).
                }\\
    
                Eres libre de compartir y redistribuir el contenido de esta obra en cualquier medio o formato, siempre y cuando des el crédito adecuado a los autores originales y no persigas fines comerciales. 
            \end{minipage}
        };
    \end{tikzpicture}
    
    
    
    % 1. Foto de fondo
    % 2. Título
    % 3. Encabezado Izquierdo
    % 4. Color de fondo
    % 5. Coord x del titulo
    % 6. Coord y del titulo
    % 7. Fecha
    % 8. Autor


}


\newcommand{\portadaBook}[8]{

    % 1. Foto de fondo
    % 2. Título
    % 3. Encabezado Izquierdo
    % 4. Color de fondo
    % 5. Coord x del titulo
    % 6. Coord y del titulo
    % 7. Fecha
    % 8. Autor

    % Personaliza el formato del título
    \pretitle{\begin{center}\bfseries\fontsize{42}{56}\selectfont}
    \posttitle{\par\end{center}\vspace{2em}}
    
    % Personaliza el formato del autor
    \preauthor{\begin{center}\Large}
    \postauthor{\par\end{center}\vfill}
    
    % Personaliza el formato de la fecha
    \predate{\begin{center}\huge}
    \postdate{\par\end{center}\vspace{2em}}
    
    \title{#2}
    \author{\href{https://losdeldgiim.github.io/}{Los Del DGIIM, \texttt{\large losdeldgiim.github.io}}
    \\ \vspace{0.5cm}#8}
    \date{Granada, #7}
    \maketitle
    
    \tableofcontents
}




\newcommand{\portadaArticle}[8]{

    % 1. Foto de fondo
    % 2. Título
    % 3. Encabezado Izquierdo
    % 4. Color de fondo
    % 5. Coord x del titulo
    % 6. Coord y del titulo
    % 7. Fecha
    % 8. Autor

    % Personaliza el formato del título
    \pretitle{\begin{center}\bfseries\fontsize{42}{56}\selectfont}
    \posttitle{\par\end{center}\vspace{2em}}
    
    % Personaliza el formato del autor
    \preauthor{\begin{center}\Large}
    \postauthor{\par\end{center}\vspace{3em}}
    
    % Personaliza el formato de la fecha
    \predate{\begin{center}\huge}
    \postdate{\par\end{center}\vspace{5em}}
    
    \title{#2}
    \author{\href{https://losdeldgiim.github.io/}{Los Del DGIIM, \texttt{\large losdeldgiim.github.io}}
    \\ \vspace{0.5cm}#8}
    \date{Granada, #7}
    \thispagestyle{empty}               % Sin encabezado ni pie de página
    \maketitle
    \vfill
}
    \portadaExamen{ffccA4.jpg}{Topología II\\Examen XI}{Topología II. Examen XI}{MidnightBlue}{-8}{28}{2026}{José Juan Urrutia Milán}

    \begin{description}
        \item[Asignatura] Topología II.
        \item[Curso Académico] 2025/26.
        \item[Grado] Doble Grado en Ingeniería Informática y Matemáticas.
        \item[Grupo] Grupo Único.
        \item[Profesor] José Antonio Gálvez.
        \item[Descripción] Examen Extraordinario.
        \item[Fecha] 6 de febrero de 2026.
        \item[Duración] 2 horas y media.
    \end{description}
    \newpage


    % ------------------------------------

    \noindent
    \textbf{Responde a cuatro de las siguientes seis preguntas.}

    \begin{ejercicio}
        Demuestra que no existe una aplicación continua e inyectiva desde la esfera $\mathbb{S}^2$ en el toro $T$.
    \end{ejercicio}

    \begin{ejercicio}
        Demuestra que el toro $T$ menos un punto no es homeomorfo al plano $\mathbb{R}^2$ menos un punto.
    \end{ejercicio}

    \begin{ejercicio}
        Consideremos la circunferencia de $\mathbb{R}^2$ dada por
        \begin{equation*}
            C_r = \{x\in \mathbb{R}^2 : \|x\| = r\}
        \end{equation*}
        con $r>0$. Si $f:\mathbb{R}^2\to \mathbb{R}^2$ es una aplicación continua tal que $f\big|_{C_r}:C_r\to \mathbb{R}^2\setminus \{(0,0)\}$ no es homotópicamente nula, demuestra que existe un punto $x_0\in \mathbb{R}^2$ tal que se cumple $f(x_0) = (0,0)$.
    \end{ejercicio}

    \begin{ejercicio}
        Consideremos la esfera $\mathbb{S}^n$ con $n\geq 2$ y $p:\mathbb{S}^n\to B$ una aplicación recubridora sobre un espacio topológico $B$. Demuestra que $\pi_1(B,b)$ es finito para cada $b\in B$.
    \end{ejercicio}

    \begin{ejercicio}
        Sean $R$ y $B$ dos espacios topológicos donde $B$ es un espacio métrico compacto. Demuestra que si $p:R\to B$ es una aplicación recubridora donde cada punto $b\in B$ tiene el mismo número finito de preimágenes $n_0\in \mathbb{N}$ entonces $R$ es compacto.
    \end{ejercicio}

    \begin{ejercicio}
        Determina la superficie compacta $S$ dada por la presentación poligonal.
        \begin{equation*}
            \langle a,b,c,d,e,f,g,h,i,j,k,l;ag^{-1}e,kb^{-1}f^{-1},h^{-1}fc,kai^{-1},jc^{-1}l^{-1},bi^{-1}g,jd^{-1}h^{-1}, edl^{-1} \rangle 
        \end{equation*}
    \end{ejercicio}

    \newpage
    \setcounter{ejercicio}{0}
    \noindent
    \textbf{Solución.}

    \begin{ejercicio}
        Demuestra que no existe una aplicación continua e inyectiva desde la esfera $\mathbb{S}^2$ en el toro $T$.\\

        \noindent
        Por reducción al absurdo, supuesto que existe una aplicación $f:\mathbb{S}^2\to T$ continua e inyectiva, consideramos una aplicación recubridora $p:\mathbb{R}^2\to T$ (que puede ser por ejemplo $p = p_0\times p_0$ donde $p_0$ es la aplicación recubridora estándar de $\mathbb{S}^1$). Fijado $x_0\in \mathbb{S}^2$, $y_0 = f(x_0)$ y tomando $r_0 \in p^{-1}(y_0)$ tenemos que $\pi_1(\mathbb{S}^2,x_0)$ es trivial, con lo que la aplicación $f$ se puede levantar a una aplicación continua $\hat{f}:\mathbb{S}^2\to \mathbb{R}^2$. Por el Teorema de Borsuk-Ulam tenemos que existe $z_0\in \mathbb{S}^2$ de forma que:
        \begin{equation*}
            \hat{f}(z_0) = \hat{f}(-z_0)
        \end{equation*}
        Por lo que tenemos que:
        \begin{equation*}
            f(z_0) = p(\hat{f}(z_0)) = p(\hat{f}(-z_0)) = f(-z_0)
        \end{equation*}
        con $z_0\neq -z_0$, hemos llegado a una contradicción con la condición de que $f$ es inyectiva.
    \end{ejercicio}

    \begin{ejercicio}
        Demuestra que el toro $T$ menos un punto no es homeomorfo al plano $\mathbb{R}^2$ menos un punto.\\

        \definecolor{torusfill}{RGB}{236, 243, 251}
        \definecolor{torusstroke}{RGB}{55, 90, 135}
        \definecolor{accentred}{RGB}{195, 35, 45}
        \definecolor{accentgreen}{RGB}{20, 135, 70}
        \definecolor{accentpurple}{RGB}{125, 45, 160}
        \definecolor{accentblue}{RGB}{25, 110, 200}

        Supongamos por reducción al absurdo que existe un homeomorfismo entre $T \setminus \{p\}$ y $\mathbb{R}^2 \setminus \{q\}$, donde $p \in T$ y $q \in \mathbb{R}^2$. Como el grupo fundamental es un invariante topológico y ambos espacios son arcoconexos, sus grupos fundamentales deberían ser isomorfos:
        \[
        \pi_1(T \setminus \{p\}) \cong \pi_1(\mathbb{R}^2 \setminus \{q\}).
        \]

        Por un lado, sabemos que $\bb{R}^2 \setminus \{q\}$ tiene por retracto de deformación a un conjunto homeomorfo a $\mathbb{S}^1$, por lo que su grupo fundamental es isomorfo a $\mathbb{Z}$. Calculemos ahora el grupo fundamental de $T \setminus \{p\}$ y veremos que no es isomorfo a $\mathbb{Z}$.\\
        

        Consideremos el toro $T \subset \mathbb{R}^3$ como una superficie de revolución obtenida al girar una circunferencia alrededor del eje vertical $Z$. Sin pérdida de generalidad, situamos el punto eliminado $p$ en el extremo derecho del ecuador exterior. Consideramos en $T$ las dos circunferencias principales que se cortan transversalmente en $p$, junto con sus circunferencias opuestas (Figura~\ref{fig:toro_cortes}):
        \begin{itemize}
            \item \textbf{Meridiano $M_p \cong \mathbb{S}^1$:} sección circular vertical que rodea el toro por la parte derecha pasando por $p$, y su \textbf{meridiano opuesto $M_{\text{izq}} \cong \mathbb{S}^1$} en la parte izquierda.
            \item \textbf{Ecuador exterior $P_p \cong \mathbb{S}^1$:} paralelo horizontal exterior del toro que pasa por $p$, y el \textbf{ecuador interior $P_{\text{int}} \cong \mathbb{S}^1$} (la circunferencia de radio mínimo que bordea el agujero central).
        \end{itemize}
        \begin{figure}[htbp]
            \centering
            \begin{tikzpicture}[scale=1.2, >=Stealth]
                % Superficie del toro
                \fill[torusfill] (0,0) ellipse (4.2 and 2.15);
                \path[fill=white]
                (-1.32, 0.02) .. controls (-0.65, -0.40) and (0.65, -0.40) .. (1.32, 0.02)
                .. controls (0.65, 0.36) and (-0.65, 0.36) .. (-1.32, 0.02);
                \draw[torusstroke, thick] (0,0) ellipse (4.2 and 2.15);
                \draw[torusstroke, thick]
                (-1.55, 0.12) .. controls (-0.8, -0.45) and (0.8, -0.45) .. (1.55, 0.12);
                \draw[torusstroke, thick]
                (-1.32, 0.02) .. controls (-0.65, 0.36) and (0.65, 0.36) .. (1.32, 0.02);

                % Ecuador exterior P_p (corte horizontal en rojo)
                \draw[accentred, very thick, dashed] (4.2, 0) arc (0:180:4.2 and 0.68);
                \draw[accentred, very thick] (-4.2, 0) arc (180:360:4.2 and 0.68);
                \node[accentred, below] at (0, -0.72) {$P_p$ (ecuador exterior)};

                % Ecuador interior P_int (borde del agujero central, en verde)
                \draw[accentgreen, line width=1.8pt, dashed]
                (1.32, 0.02) .. controls (0.65, 0.36) and (-0.65, 0.36) .. (-1.32, 0.02);
                \draw[accentgreen, line width=1.8pt]
                (-1.32, 0.02) .. controls (-0.65, -0.40) and (0.65, -0.40) .. (1.32, 0.02);
                \node[accentgreen, above] at (0, 0.72) {$P_{\text{int}}$ (ecuador interior)};

                % Meridiano M_p pasando por p (corte vertical derecho en rojo)
                \draw[accentred, very thick]
                (1.32, 0.02) .. controls (1.75, 0.72) and (3.85, 0.72) .. (4.2, 0);
                \draw[accentred, very thick, dashed]
                (4.2, 0) .. controls (3.85, -0.68) and (1.75, -0.68) .. (1.32, 0.02);
                \node[accentred, above right] at (2.85, 0.68) {$M_p$ (meridiano derecho)};

                % Meridiano opuesto M_op (retracto morado en el tubo izquierdo)
                \draw[accentpurple, very thick]
                (-4.2, 0) .. controls (-3.85, 0.72) and (-1.75, 0.72) .. (-1.32, 0.02);
                \draw[accentpurple, very thick, dashed]
                (-1.32, 0.02) .. controls (-1.75, -0.68) and (-3.85, -0.68) .. (-4.2, 0);
                \node[accentpurple, above left] at (-2.85, 0.68) {$M_{\text{izq}}$ (meridiano izquierdo)};

                % Punto p = M_p \cap P_p
                \fill[white] (4.2, 0) circle (2.8pt);
                \draw[accentred, very thick] (4.2, 0) circle (2.8pt)
                node[right=5pt, accentred] {$p = M_p \cap P_p$};
            \end{tikzpicture}
            \caption{Cortes sobre el toro $T$}
            \label{fig:toro_cortes}
        \end{figure}

        Observemos que $M_p \cap P_p = \{p\}$. Definimos los subconjuntos de $X = T \setminus \{p\}$:
        \[
        U = X \setminus M_p = T \setminus M_p, \qquad V = X \setminus P_p = T \setminus P_p.
        \]

        Observamos que $X=U \cup V$ y que $U,V$ son abiertos en $X$.
        \begin{itemize}
            \item Estudiamos $U$. Al eliminar el meridiano $M_p$, cortamos transversalmente el tubo del toro, por lo que $U$ tiene como retracto de deformación al meridiano opuesto $M_{\text{izq}} \cong \mathbb{S}^1$. Por tanto, $U$ es arcoconexo y tiene grupo fundamental isomorfo a $\mathbb{Z}$.
            \item Estudiamos $V$. Al eliminar el ecuador exterior $P_p$, abrimos el toro a lo largo de su circunferencia exterior de radio máximo, de modo que $V$ tiene como retracto de deformación al ecuador interior $P_{\text{int}} \cong \mathbb{S}^1$. Por tanto, $V$ es arcoconexo y tiene grupo fundamental isomorfo a $\mathbb{Z}$.
            \item Estudiamos $U \cap V$. Aplicando primero el retracto de deformación de $U$ sobre $M_{\text{izq}}$ y luego el retracto de deformación de $V$ sobre $P_{\text{int}}$, obtenemos que $U \cap V$ tiene como retracto de deformación un único punto, por lo que es contráctil, luego simplemente conexo.
        \end{itemize}

        Aplicando el Teorema de Seifert-van Kampen (en el caso en que la intersección tiene grupo fundamental trivial), obtenemos:
        \begin{equation*}
            \pi_1(T \setminus \{p\}) \cong \pi_1(U) \ast \pi_1(V) \cong \mathbb{Z} \ast \mathbb{Z}.
        \end{equation*}
        Además, puesto que es la unión de dos subespacios arcoconexos con intersección no vacía, $T \setminus \{p\}$ es arcoconexo, y por tanto $\pi_1(T \setminus \{p\})$ no depende del punto base elegido.\\

        De esta forma, puesto que $\bb{Z}$ es un grupo abeliano y $\bb{Z} \ast \bb{Z}$ no lo es, concluimos que $\pi_1(T \setminus \{p\})$ y $\pi_1(\mathbb{R}^2 \setminus \{q\})$ no son isomorfos, y por tanto $T \setminus \{p\}$ y $\mathbb{R}^2 \setminus \{q\}$ no son homeomorfos.
    \end{ejercicio}

    \begin{ejercicio}
        Consideremos la circunferencia de $\mathbb{R}^2$ dada por
        \begin{equation*}
            C_r = \{x\in \mathbb{R}^2 : \|x\| = r\}
        \end{equation*}
        con $r>0$. Si $f:\mathbb{R}^2\to \mathbb{R}^2$ es una aplicación continua tal que $f\big|_{C_r}:C_r\to \mathbb{R}^2\setminus \{(0,0)\}$ no es homotópicamente nula, demuestra que existe un punto $x_0\in \mathbb{R}^2$ tal que se cumple $f(x_0) = (0,0)$.\\

        Supongamos por reducción al absurdo que $f(x) \neq (0,0)$ para todo $x \in \mathbb{R}^2$. Entonces podemos considerar la siguiente restricción en el codominio, que sigue siendo continua:
        \begin{equation*}
            f:\mathbb{R}^2 \to \mathbb{R}^2 \setminus \{(0,0)\}
        \end{equation*}

        Puesto que $\bb{R}^2$ es contráctil desde el punto $(0,0)$, tenemos la siguiente homotopía:
        \Func{H}{\bb{R}^2 \times [0,1]}{\bb{R}^2}{(x,t)}{(1-t)x}

        Por tanto, consideramos la siguiente aplicación continua:
        \Func{g}{C_r \times [0,1]}{\bb{R}^2 \setminus \{(0,0)\}}{(x,t)}{f(H(x,t)) = f((1-t)x)}

        Esta aplicación es continua y está bien definida, pues $f$ no toma el valor $(0,0)$ en ningún punto. Además, tenemos que:
        \begin{itemize}
            \item $g(x,0) = f(H(x,0)) = f(x)$
            \item $g(x,1) = f(H(x,1)) = f(0,0)$
        \end{itemize}

        Por tanto, $g$ es una homotopía entre $f\big|_{C_r}$ y la aplicación constante $f(0,0)$, lo que contradice la hipótesis de que $f\big|_{C_r}$ no es homotópicamente nula. Esta contradicción viene de suponer que $f(x) \neq (0,0)$ para todo $x \in \mathbb{R}^2$, por lo que debe existir al menos un punto $x_0 \in \mathbb{R}^2$ tal que $f(x_0) = (0,0)$.
    \end{ejercicio}

    \begin{ejercicio}
        Consideremos la esfera $\mathbb{S}^n$ con $n\geq 2$ y $p:\mathbb{S}^n\to B$ una aplicación recubridora sobre un espacio topológico $B$. Demuestra que $\pi_1(B,b)$ es finito para cada $b\in B$.\\

        Sea $b \in B$ y sea $x \in p^{-1}(b)$. Como $\bb{S}^n$ es simplemente conexo para $n \geq 2$, tenemos que la correspondencia del levantamiento es biyectiva, luego consideramos la siguiente biyección:
        \Func{\phi}{\pi_1(B,b)}{p^{-1}(b)}{[\gamma]}{\tilde{\gamma}_x(1)}

        Por tanto, de ser $p^{-1}(b)$ finito, $\pi_1(B,b)$ también lo sería. Como $\bb{S}^n$ es compacto, en el Examen VI se demostró que $p^{-1}(b)$ es finito, y por tanto $\pi_1(B,b)$ también lo es.
    \end{ejercicio}

    \begin{ejercicio}
        Sean $R$ y $B$ dos espacios topológicos donde $B$ es un espacio métrico compacto. Demuestra que si $p:R\to B$ es una aplicación recubridora donde cada punto $b\in B$ tiene el mismo número finito de preimágenes $n_0\in \mathbb{N}$ entonces $R$ es compacto.
    \end{ejercicio}

    \begin{ejercicio}
        Determina la superficie compacta $S$ dada por la presentación poligonal.
        \begin{equation*}
            \langle a,b,c,d,e,f,g,h,i,j,k,l;ag^{-1}e,kb^{-1}f^{-1},h^{-1}fc,kai^{-1},jc^{-1}l^{-1},bi^{-1}g,jd^{-1}h^{-1}, edl^{-1} \rangle 
        \end{equation*}
    \end{ejercicio}

\end{document}
